\documentclass[11pt]{article}
\usepackage[T1]{fontenc}
\usepackage{amsmath,amssymb,amsthm}
\usepackage[margin=1in]{geometry}
\usepackage{booktabs}
\usepackage[hidelinks,pdfusetitle]{hyperref}

\newtheorem{theorem}{Theorem}
\newtheorem{lemma}[theorem]{Lemma}

\title{A Complete Margin-One Five-Voter Census\\
for 254 Nine-Vertex Boundary Tournaments}
\author{Xiaojian Zeng\\
PBC School of Finance, Tsinghua University, 100190\\
\texttt{xiaojianzeng@foxmail.com}}
\date{}

\begin{document}
\maketitle

\begin{abstract}
Chindelevitch and Harutyunyan reported 254 non-isomorphic tournaments on nine
vertices that are inducible by three voters but not by a three-voter profile
of margin one on every arc.  We determine the corresponding finite question
at five voters: every one of these 254 tournaments has a five-voter profile in
which every majority arc is supported by exactly three voters.  The proof is a
complete collection of 1,270 explicit permutations, checked on all 9,144
tournament arcs by a separate exact verifier.  Thus the reported boundary
family is contained in $I_{5,1}$.  This finite census does not settle whether
the global inclusion $I_{5,1}\subseteq I_{5,3}$ is strict.
\end{abstract}

\section{Statement and context}

Let $T=(V,A)$ be a tournament.  For an odd integer $m$, an $m$-profile
$P=(\pi_1,\ldots,\pi_m)$ is a list of strict total orders of $V$; repetitions
are allowed.  If $x\to y$ is an arc of $T$, write
\[
 s_P(x,y)=\bigl|\{r:x\text{ precedes }y\text{ in }\pi_r\}\bigr|,
 \qquad \mu_P(x,y)=2s_P(x,y)-m.
\]
For odd $t\leq m$, let $I_{m,t}$ be the class of tournaments admitting an
$m$-profile for which $s_P(x,y)>(m/2)$ and $\mu_P(x,y)\leq t$ on every arc.
These margin classes and the motivating hierarchy are defined in
\cite{ChindelevitchHarutyunyan2026}.  In particular, a certificate for
$T\in I_{5,1}$ consists of five orders with $s_P(x,y)=3$ for every arc.

The same source reports exactly 254 non-isomorphic nine-vertex tournaments in
$I_{3,3}\setminus I_{3,1}$ and supplies their serialized adjacency matrices
\cite{ChindelevitchHarutyunyan2026,CHData2026}.  It also conjectures strictness
of every adjacent margin layer.  At five voters, the first inclusion in that
conjecture is
\[
                 I_{5,1}\subsetneq I_{5,3}.
\]
The 254 tournaments form a natural first boundary family to test because each
lies in $I_{5,3}$: append an arbitrary order and its reverse to an inducing
three-profile, thereby leaving every majority margin unchanged.

\begin{theorem}\label{thm:census}
Every one of the 254 source-reported non-isomorphic nine-vertex tournaments
in $I_{3,3}\setminus I_{3,1}$ belongs to $I_{5,1}$.
\end{theorem}

The theorem resolves this finite question positively.  It neither proves nor
disproves the displayed global strict inclusion.

\section{Certificate and proof}

Label the vertices $0,\ldots,8$.  For each representative, the accompanying
certificate gives a 36-bit word in lexicographic upper-triangle order:
for $i<j$, bit $1$ denotes $i\to j$ and bit $0$ denotes $j\to i$.  The same
record gives five permutations of $0,\ldots,8$, read from first to last.
All 254 adjacency words and all five permutations for every word are contained
in the machine-readable supplement
\texttt{evidence/research\_pass2\_results.json}.

For orientation, the first record has adjacency word
\[
\texttt{000001001100000001010110000001110101}
\]
and the following five orders:
\begin{align*}
 &(8,6,5,4,1,3,2,7,0),\\
 &(1,5,0,6,2,7,3,4,8),\\
 &(8,7,1,0,3,4,2,5,6),\\
 &(6,2,3,5,4,0,7,1,8),\\
 &(7,4,8,3,2,0,5,6,1).
\end{align*}
This displayed case is only an illustration; the theorem uses every record in
the supplement.

\begin{lemma}\label{lem:local}
Suppose a certificate record contains a tournament adjacency word and five
permutations such that, for each of the 36 unordered vertex pairs, exactly
three permutations place the tail of the encoded arc before its head.  Then
the encoded tournament belongs to $I_{5,1}$.
\end{lemma}

\begin{proof}
Each permutation is a strict total order.  On an encoded arc $x\to y$, the
hypothesis gives $s_P(x,y)=3$, so the five orders induce that arc and its
margin is $2\cdot3-5=1$.  This holds for every pair, hence the profile witnesses
membership in $I_{5,1}$.
\end{proof}

\begin{proof}[Proof of Theorem~\ref{thm:census}]
The standalone verifier
\texttt{evidence/research\_pass2\_verify.py} reads the certificate JSON and
checks, for each record, that the adjacency word has 36 binary entries and
that each of the five listed orders is a permutation of $0,\ldots,8$.  It then
constructs inverse positions and recounts support on all 36 pairs.  Every
encoded target arc has support exactly three.  Applying
Lemma~\ref{lem:local} to all 254 records proves that every serialized
representative lies in $I_{5,1}$.

The verifier also checks the adjacency digests and confirms by exact directed
graph isomorphism tests that the 254 representatives are pairwise
non-isomorphic.  Membership in $I_{5,1}$ is invariant under relabelling, so one
certificate for each representative proves the asserted classwise statement.
\end{proof}

The verification totals are recorded in Table~\ref{tab:totals}.  They are
counts of exact comparisons, not floating-point tests.

\begin{table}[ht]
\centering
\begin{tabular}{lr}
\toprule
Checked object & Count \\
\midrule
Tournament representatives & 254 \\
Five-order profiles & 254 \\
Individual permutations & 1,270 \\
Arc-support recounts & 9,144 \\
Degree-compatible isomorphism pairs tested & 5,284 \\
\bottomrule
\end{tabular}
\caption{Exhaustive checks performed by the standalone verifier.}
\label{tab:totals}
\end{table}

The certificate architecture deliberately separates profile discovery from
the proof check: solver output is not trusted after the explicit permutations
have been emitted.  A related separation between generation and independent
exact verification is used for finite discrete models in
\cite{ZengEtAl2026FiniteCertificates}; no mathematical result from that paper
is needed here.

\section{Reproducibility and limits}

The source data are preserved as
\texttt{evidence/n9\_margin1only\_tournaments.rds.base64}.  The program
\texttt{evidence/research\_pass2.py} parses the complete version-3 XDR R
serialization, validates every $9\times9$ matrix, and emits the certificate
records.  It used an SMT solver only to find the orders.  The proof-bearing
verification command is
\begin{verbatim}
python3 evidence/research_pass2_verify.py
\end{verbatim}
Its saved output is
\texttt{evidence/research\_pass2\_verification.json}.  The verifier imports
neither the XDR parser nor the solver code and directly checks the emitted
finite witnesses.  An independent audit additionally reparsed the pinned RDS
bytes and matched all 254 adjacency words before repeating the full support
recount.

The source paper and data repository identify the serialized list as the
reported boundary family.  The present computation checks the list's size,
tournament identities, pairwise non-isomorphism, and all claimed five-order
profiles; it does not redo the substantially larger source enumeration that
selected these 254 classes.  There are no negative membership instances in
this census, so no unsatisfiability certificate is involved.  Finally, this
result supplies no witness in $I_{5,3}\setminus I_{5,1}$ and makes no claim
about the minimum order of a tournament not inducible by five voters.

\paragraph{Computational and AI assistance.}
OpenAI Codex assisted with literature search, profile-search code, exact
verification code, audit organization, and initial manuscript preparation.
Z3 version 4.15.3 was used to discover the five-order profiles.  The theorem
does not rely on a solver status: all profiles are explicit and their defining
property is recounted by a separate checker.  No proof assistant was used.

\bibliographystyle{plain}
\bibliography{references}

\end{document}
